% Gerado por artigo/analise/para_latex.py a partir de artigo/resultados/numeros.json. Não editar à mão.
\begin{table}[htbp]
\centering
\fontsize{9}{11}\selectfont
\setlength{\tabcolsep}{4pt}
\caption{Conduta registrada ao fim dos atendimentos na atenção primária à saúde, por diagnóstico. Brasil, janeiro de 2025 a set/2026.}
\label{tab:tabela3}
\begin{tabular}{>{\raggedright\arraybackslash}p{2.5cm}R{1.9cm}R{1.4cm}R{2.0cm}R{2.0cm}R{2.5cm}R{2.4cm}}
\toprule
Diagnóstico & Atendimentos (n) & Alta do episódio (\%) & Retorno para consulta agendada (\%) & Retorno para cuidado continuado (\%) & Encaminhamento para serviço especializado (\%) & Encaminhamento para CAPS (\%) \\ \midrule
Jogo patológico (F63.0) & 8.919 & 37,6 & 35,2 & 31,4 & 24,2 & 1,9 \\
Ansiedade (F41) & 2.387.047 & 48,6 & 27,9 & 28,7 & 9,9 & 0,2 \\
Depressão (F32) & 753.674 & 48,6 & 27,4 & 29,5 & 11,6 & 0,6 \\
Uso de álcool (F10) & 66.420 & 54,3 & 24,4 & 25,4 & 11,3 & 1,8 \\
\bottomrule
\end{tabular}
\par\smallskip\parbox{\linewidth}{\raggedright CAPS: Centro de Atenção Psicossocial. Um atendimento pode ter mais de uma conduta, e as linhas não somam 100\%. O registro não identifica o serviço especializado. Fonte: Sistema de Informação em Saúde para a Atenção Básica.}
\end{table}
